\documentclass[10pt,a4paper]{amsart}
\usepackage[margin=19mm]{geometry}
\usepackage{amsmath,amssymb,mathtools}
\usepackage[scr=rsfs,cal=euler]{mathalfa}
\usepackage{xfrac}
\usepackage[unicode,hidelinks,bookmarks=false]{hyperref}
\newcommand{\ie}{i.e.\ }
\newcommand{\cf}{cf.\ }
\newcommand{\resp}{respectively\ }
\newcommand{\Subsection}{\subsection}
\providecommand{\nobreakdash}{\nobreak\hbox{-}}
\setlength{\emergencystretch}{2em}
\multlinegap0pt
\usepackage{float}
\usepackage{tikz}
\hypersetup{urlcolor=blue, citecolor=red}
     \setcounter{page}{1}
\newtheorem{theorem}{Theorem}[section]
\newtheorem{corollary}[theorem]{Corollary}
\newtheorem*{main}{Main Theorem}
\newtheorem{lemma}[theorem]{Lemma}
\newtheorem{proposition}[theorem]{Proposition}
\newtheorem{conjecture}{Conjecture}
\newtheorem{assumption}[theorem]{Assumption}
\newtheorem*{problem}{Problem}
\newtheorem{definition}[theorem]{Definition}
\newtheorem{remark}[theorem]{Remark}
\newtheorem*{notation}{Notation}
\newcommand{\ep}{\varepsilon}
\newcommand{\R}{\mathbb{R}}
\newcommand{\N}{\mathbb{N}}
\newcommand{\C}{\mathbb{C}}
\newcommand{\T}{\mathbb{T}}
\newcommand{\Z}{\mathbb{Z}}
\newcommand{\D}{\mathbb{D}}
\newcommand{\K}{\mathbb{K}}
\newcommand{\Q}{\mathbb{Q}}
\newcommand{\bl}{\textcolor{blue}}
\newcommand{\red}{\textcolor{red}}
\begin{document}
\title[Low-regularity well-posedness for dispersive equations with ]{Low-regularity well-posedness for dispersive equations with derivative nonlinearity and quasi-periodic initial data}
\author{Robert Schippa}
\date{}
\maketitle
\noindent\emph{Source and licence.} Low-regularity well-posedness for dispersive equations with derivative nonlinearity and quasi-periodic initial data. Version arXiv:2608.09470v1 (CC BY 4.0). This material is distributed under the Creative Commons Attribution 4.0 International licence, http://creativecommons.org/licenses/by/4.0/, as recorded in the arXiv OAI licence metadata for this exact version and shown on the abstract page https://arxiv.org/abs/2608.09470v1. Full rights notice: licenses/arxiv-2608.09470-CC-BY-4.0.txt.\par\smallskip
\noindent\textit{Editorial scope.} Counted unit: the original \texttt{theorem} environment \texttt{thm:SFBO} at source lines 499--510 together with its complete proof at lines 511--541; the delivered document keeps source lines 402--542 so that every referenced label and every citation resolves inside it. Native editable source is the paper's own concatenated TeX; the AMS wrapper is editorial formatting only. No publisher class, upstream script or unrelated artwork is redistributed.
\section{Multilinear estimates in adapted function spaces}
\label{section:Trilinear}
In this section we prove bi- and trilinear estimates in function spaces adapted to dispersive equations. Together with the frequency-dependent time localization, these will allow us to control the derivative nonlinearity.
We begin with a bilinear estimate, which relies on the bilinear C\'ordoba--Fefferman square function estimate. First, we recall the linear C\'ordoba--Fefferman square function estimate. Let $\pm \mathbb{P}^1 = \{ (\xi,\pm \xi^2) : 0 < \xi < 2 \}$ denote the graph of the positive resp. negative parabola, and $\mathcal{N}_{\delta}(\pm \mathbb{P}^1)$ denote the $\delta$-neighborhood of $\pm \mathbb{P}^1$. For $0< \alpha \leq 1$ we denote by $\mathbb{I}_{\alpha}$ a finitely overlapping partition of $(-3,3)$ into $\alpha$-intervals.

For $F \in \mathcal{S}(\R^2)$, $\theta \in \mathbb{I}_{\alpha}$ we denote by $F_{\theta}$ the Fourier projection  of the first frequency coordinate to $\theta$ given by $\hat{F}_{\theta}(\xi) = 1_{\theta}(\xi_1) \hat{F}(\xi_1,\xi_2)$. By $\pi_1 : \R^2 \to \R$ the projection $(x_1,x_2) \mapsto x_1$ is deonted. We have the following square function estimate:
\begin{theorem}[C\'ordoba--Fefferman square function estimate]
Let $0<\delta \leq 1$, $F \in \mathcal{S}(\R^2)$ with $\text{supp}(\hat{F}) \subseteq \mathcal{N}_{\delta}(\pm \mathbb{P}^1)$. Then the following holds:
\begin{equation*}
\| F \|_{L^4(\R^2)} \lesssim \big\| \big( \sum_{\theta \in \mathbb{I}_{\delta^{1/2}}} |F_{\theta}|^2 \big)^{\frac{1}{2}} \big\|_{L^4(\R^2)}.
\end{equation*}
\end{theorem}
The square function estimate depends on the non-degeneracy of the curve and versions for $\Gamma_3 = \{ (\xi,\xi^3) : \xi \in (-2,2) \}$ need to respect the flatness at the origin. In \cite[Theorem~2.2]{Schippa2025Quasiperiodic} a square function estimate with $\delta^{1/3}$-intervals is proved.

The square function estimate can be used to estimate exponential sums by approximating the exponential sum with an oscillatory integral, carrying out the square function estimate and passing back to an exponential sum. This will be referred to as continuous approximation in the following. To the best of the author's knowledge, this approximation was first carried out by Bourgain \cite{Bourgain2013}.
In \cite[Section~3.1]{Schippa2025Quasiperiodic} was pointed out how the square function estimate together with the continuous approximation yields sharp estimates for quasi-periodic functions. 

\smallskip

Here we point out how the bilinear C\'ordoba--Fefferman square function estimate yields new bilinear estimates for exponential sums related to quasi-periodic functions. Under an additional transversality assumption we have the following refinement of the linear estimate:
\begin{theorem}[Bilinear C\'ordoba--Fefferman square function estimate for the cubic]
Let $0 < \delta \leq 1$, $F_i \in \mathcal{S}(\R^2)$, $i=1,2$ with $\text{supp}(\hat{F}_i) \subseteq \mathcal{N}_{\delta}(\Gamma_3)$ such that
\begin{equation*}
\text{dist}(\text{supp}(\pi_1 \hat{F}_1),\text{supp}(\pi_1 \hat{F}_2)) \geq 1/16 \text{ and }
\text{dist}(\text{supp}(\pi_1 \hat{F}_1),\text{supp}(- \pi_1 \hat{F}_2)) \geq 1/16.
\end{equation*}
Then it holds:
\begin{equation*}
\| F_1 F_2 \|_{L^2(\R^2)} \lesssim \big\| \big( \sum_{\theta_1 \in \mathbb{I}_{\delta}} |F_{\theta_1}|^2 \big)^{\frac{1}{2}} \big( \sum_{\theta_2 \in \mathbb{I}_{\delta}} |F_{\theta_2}|^2 \big)^{\frac{1}{2}} \big\|_{L^2(\R^2)}.
\end{equation*}
\end{theorem}
\begin{proof}
We can suppose by Cauchy-Schwarz that $\delta < 1/32$.
We let $F_i = \sum_{\theta \in \mathbb{I}_{\delta}} F_{i \theta}$ and obtain
\begin{equation*}
\big\| \sum_{\theta_1 \in \mathbb{I}_{\delta}} F_{1 \theta_1} \sum_{\theta_2 \in \mathbb{I}_{\delta}} F_{2 \theta_2} \big\|^2_{L^2(\R^2)} = \sum_{\substack{\theta_1,\theta_2, \\ \theta_3,\theta_4}} \int F_{1 \theta_1} F_{2 \theta_2} \overline{F}_{1 \theta_3} \overline{F}_{2 \theta_4} dx.
\end{equation*}
We find by Plancherel's theorem that
\begin{equation}
\label{eq:TrivialBiorthogonality}
\int F_{1 \theta_1} F_{2 \theta_2} \overline{F}_{1 \theta_3} \overline{F}_{2 \theta_4} dx = 0,
\end{equation}
unless there are $\xi_i \in 2 \theta_i$, which satisfy
\begin{equation*}
\left\{ \begin{array}{cl}
\xi_1 + \xi_2 &= \xi_3 + \xi_4, \\
\xi_1^3 + \xi_2^3 &= \xi_3^3 + \xi_4^3 + \mathcal{O}(\delta).
\end{array} \right.
\end{equation*}
Taking the third power of the first equation and subtracting the second, we find
\begin{equation*}
\left\{ \begin{array}{cl}
\xi_1 + \xi_2 &= \xi_3 + \xi_4, \\
\xi_1 \xi_2 (\xi_1+\xi_2) &= \xi_3 \xi_4(\xi_3 + \xi_4) + \mathcal{O}(\delta).
\end{array} \right.
\end{equation*}
By the separation of $\pi_1 \text{supp}(\hat{F}_1)$ and $-\pi_1 \text{supp}(\hat{F}_2)$, we obtain
\begin{equation*}
\left\{ \begin{array}{cl}
\xi_1 + \xi_2 &= \xi_3 + \xi_4, \\
2 \xi_1 \xi_2  &=  2 \xi_3 \xi_4 + \mathcal{O}(\delta).
\end{array} \right.
\end{equation*}
This is equivalent to
\begin{equation}
\label{eq:BiorthogonalityAux}
\left\{ \begin{array}{cl}
\xi_1 + \xi_2 &= \xi_3 + \xi_4, \\
\xi_1^2 + \xi_2^2  &=  \xi_3^2 + \xi_4^2 + \mathcal{O}(\delta)
\end{array} \right.
\end{equation}
and by the classical bilinear C\'ordoba--Fefferman argument, which depends on \\ $\text{dist}(\pi_1 \text{supp}(\hat{F}_1),\pi_1 \text{supp}(\hat{F}_2)) \geq 1/16$, we find that \eqref{eq:TrivialBiorthogonality} holds unless
\begin{equation}
\label{eq:Biorthogonality}
|\xi_1 - \xi_3| \lesssim \delta, \quad |\xi_2 - \xi_4 | \lesssim \delta,
\end{equation}
from which the claim is immediate by Cauchy-Schwarz. 

For self-containedness we note that \eqref{eq:BiorthogonalityAux} yields indeed
\begin{equation*}
\left\{ \begin{array}{cl}
\xi_1 + \xi_2 &= \xi_3 + \xi_4, \\
(\xi_1- \xi_4)(\xi_1+\xi_4)  &=  (\xi_3-\xi_2)(\xi_3+\xi_2) + \mathcal{O}(\delta)
\end{array} \right.
\end{equation*}
and furthermore, by the separation condition,
\begin{equation*}
\left\{ \begin{array}{cl}
\xi_1 + \xi_2 &= \xi_3 + \xi_4, \\
\xi_1+\xi_4  &=  \xi_3+\xi_2 + \mathcal{O}(\delta),
\end{array} \right.
\end{equation*}
which yields again \eqref{eq:Biorthogonality}.
\end{proof}

The argument is flexible and foremostly hinges on adequate separation conditions of the frequencies. We record the following variant, which plays the key role when estimating solutions to the Benjamin-Ono equation.


\begin{theorem}[C\'ordoba--Fefferman bilinear square function estimate, II]
\label{thm:SFBO}
Let $0<\delta\leq 1$, $F_i \in \mathcal{S}(\R^2)$, $i=1,2$ with $\text{supp}(\hat{F}_i) \subseteq \mathcal{N}_{\delta}(\mathbb{P}^1)$ or $\text{supp}(\hat{F}_i) \subseteq \mathcal{N}_{\delta}(-\mathbb{P}^1)$ and 
\begin{equation*}
\text{dist}(\text{supp}(\pi_1 \hat{F}_1),\text{supp}(\pi_1 \hat{F}_2)) \geq 1/16 \text{ and }
\text{dist}(\text{supp}(\pi_1 \hat{F}_1),\text{supp}(- \pi_1 \hat{F}_2)) \geq 1/16.
\end{equation*}
Then it holds:
\begin{equation*}
\| F_1 F_2 \|_{L^2(\R^2)} \lesssim \big\| \big( \sum_{\theta_1 \in \mathbb{I}_{\delta}} |F_{\theta_1}|^2 \big)^{\frac{1}{2}} \big( \sum_{\theta_2 \in \mathbb{I}_{\delta}} |F_{\theta_2}|^2 \big)^{\frac{1}{2}} \big\|_{L^2(\R^2)}.
\end{equation*}
\end{theorem}

\begin{proof}
By the same initial steps like above, we find
\begin{equation*}
\int F_{1 \theta_1} F_{2 \theta_2} \overline{F}_{1 \theta_3} \overline{F}_{2 \theta_4} = 0
\end{equation*}
unless for $\xi_i \in \theta_i$ it holds
\begin{equation*}
\left\{ \begin{array}{cl}
\xi_1 + \xi_2 &= \xi_3 + \xi_4, \\
a_1 \xi_1^2 + a_2 \xi_2^2 &= a_1 \xi_3^2 + a_2 \xi_4^2 + \mathcal{O}(\delta)
\end{array} \right.
\end{equation*}
with
\begin{equation*}
a_i = \begin{cases}
1, &\quad \text{supp}(\hat{F}_i) \subseteq \mathcal{N}_{\delta}(\mathbb{P}^1), \\
-1, &\quad \text{supp}(\hat{F}_i) \subseteq \mathcal{N}_{\delta}(-\mathbb{P}^1).
\end{cases}
\end{equation*}
By symmetry it suffices to consider $(a_1,a_2) = (1,1)$ and $(a_1,a_2) = (1,-1)$. The first case was covered above.
%We turn to the first case. in Since $|\xi_1-\xi_2| \geq 1/8$, rewriting the second equation in \eqref{eq:Biorthogonality} yields
%\begin{equation*}
%(\xi_1-\xi_4)(\xi_1+\xi_4) = (\xi_3-\xi_2)(\xi_3+\xi_2) + \mathcal{O}(\delta) \Rightarrow \xi_1 + \xi_4 = \xi_2 + \xi_3 + \mathcal{O}(\delta).
%\end{equation*}
%Hence, $\text{dist}(\theta_1,\theta_3) + \text{dist}(\theta_2,\theta_4) \leq C \delta$ for which reason the claim follows from Cauchy-Schwarz.
When $(a_1,a_2)=(1,-1)$, rewriting the second equation and the separation of the Fourier support projected to the first coordinate gives
\begin{equation*}
(\xi_1-\xi_2)(\xi_1+\xi_2) = (\xi_3-\xi_4)(\xi_3+\xi_4) + \mathcal{O}(\delta) \Rightarrow \xi_1 - \xi_2 = \xi_3 - \xi_4 + \mathcal{O}(\delta).
\end{equation*}
Hence, $\text{dist}(\theta_1,\theta_3) + \text{dist}(\theta_2,\theta_4) \leq C \delta$ from which the claim is immediate.
\end{proof}


\begin{thebibliography}{99}
\bibitem[Bourgain2013]{Bourgain2013} J.~Bourgain. \newblock Moment inequalities for trigonometric polynomials with spectrum in curved hypersurfaces. \newblock {\em Israel J. Math.}, 193(1):441--458, 2013.
\bibitem[Schippa2025Quasiperiodic]{Schippa2025Quasiperiodic} Robert Schippa. \newblock Strichartz estimates for quasi-periodic functions and applications. \newblock {\em Proc. Lond. Math. Soc. (3)}, 130(6):Paper No. e70058, 38, 2025.
\end{thebibliography}
\end{document}
